% !PRINTED 58
\weberpage{58}
%<<CONT

woraus sich ergiebt, dass in (4) irgend ein Product $x_1^{\alpha_1}x_2^{\alpha_2}\cdots$ genau
\[
 {\Pi(n)\over \Pi(\alpha_1)\Pi(\alpha_2)\cdots}
\]
mal vorkommt. Setzt man also noch fest, dass $A_{\nu_1,\nu_2,\ldots,\nu_n}$ sich nicht ändern soll, wenn die Indices beliebig permutirt werden, so erweisen sich die Bezeichnungsweisen (2) und (4) als identisch, wenn durch Zusammenfassen gleicher Factoren
\[
 x_{\nu_1}x_{\nu_2}\cdots x_{\nu_n}=x_1^{\alpha_1}x_2^{\alpha_2}\cdots x_m^{\alpha_m}
\]
und
\[
 A_{\nu_1,\nu_2,\ldots,\nu_n}=A_{\alpha_1,\alpha_2,\ldots,\alpha_m}
\]
gesetzt wird.

Bezeichnen wir die Anzahl der Glieder, die in der Function $\Phi$ [nach (2)] auftreten, mit $(m,n)$, so findet man, indem man zunächst die Glieder zählt, die den Factor $x_1$ haben und dann die übrigen, die eine homogene Function $n$ter Ordnung von den übrigen $m-1$ Variablen bilden, die Recursionsformel
\begin{equation}
 (m,n)=(m,n-1)+(m-1,n),
\tag{5}
\end{equation}
mit deren Hülfe man durch vollständige Induction den Ausdruck
\begin{equation}
 (m,n)={m(m+1)\cdots(m+n-1)\over 1\cdot2\cdots n}
 ={ \Pi(m+n-1)\over \Pi(n)\Pi(m-1)}
\tag{6}
\end{equation}
als richtig erweist.

Die ganzen homogenen Functionen werden auch \gesperrt{Formen} genannt. Man unterscheidet nach der Anzahl der Variablen \gesperrt{unäre} (einfache Potenzen), \gesperrt{binäre, ternäre, quaternäre Formen}. Die binären Formen sind es, die uns hier besonders interessiren, deren Theorie im Wesentlichen identisch ist mit der Theorie der ganzen rationalen Functionen einer Veränderlichen. Man gelangt von den binären Formen zu diesen Functionen zurück, wenn man eine der homogenen Variablen als constant ansieht, z. B. ihr den Werth $1$ giebt.
